% Source: https://artofproblemsolving.com/community/c3h1500195_please_help
% Stable identity: aops-topic-1500195-q01
\textbf{Upper bound.} If two candles burn for $T$ hours, the total burning time used is $2T$ candle-hours. The candles provide only $3+4+5=12$ candle-hours, so $2T\le12$ and $T\le6$.

\textbf{A schedule for $6$ hours.} Call the candles $A$ ($3$ h), $B$ ($4$ h) and $C$ ($5$ h). Measure time in hours from the start:
\[
\begin{array}{c|c}
\text{time} & \text{burning}\\\hline
0\text{ to }1 & B,\ C\\
1\text{ to }3 & A,\ C\\
3\text{ to }5 & B,\ C\\
5\text{ to }6 & A,\ B
\end{array}
\]
Candle $A$ burns $2+1=3$ hours, $B$ burns $1+2+1=4$ hours and $C$ burns $1+2+2=5$ hours, and exactly two candles burn at every moment.

So the greatest time is $\boxed{6}$ hours.
